%------------------------
% Ethan's Résumé Template
% Author: necusjz
% License: CC-BY-4.0
%------------------------
\documentclass[letterpaper,11pt]{article}

\usepackage{latexsym}
\usepackage[empty]{fullpage}
\usepackage{titlesec}
\usepackage{marvosym}
\usepackage[usenames,dvipsnames]{color}
\usepackage{verbatim}
\usepackage{enumitem}
\usepackage[hidelinks]{hyperref}
\usepackage{fancyhdr}
\usepackage[english]{babel}
\usepackage{tabularx}
\usepackage{fontawesome5}
\usepackage{ragged2e}
\usepackage{etoolbox}
\usepackage{tikz}
\usepackage{kotex}
\input{glyphtounicode}

% font options
\usepackage{newpxtext}
\linespread{1.05}  % palladio needs more leading (space between lines)
\usepackage[T1]{fontenc}

\pagestyle{fancy}
\fancyhf{}  % clear all header and footer fields
\fancyfoot{}
\renewcommand{\headrulewidth}{0pt}
\renewcommand{\footrulewidth}{0pt}

% adjust margins
\addtolength{\oddsidemargin}{-0.5in}
\addtolength{\evensidemargin}{-0.5in}
\addtolength{\textwidth}{1in}
\addtolength{\topmargin}{-.5in}
\addtolength{\textheight}{1.0in}

\urlstyle{same}

\raggedbottom
\raggedright
\setlength{\tabcolsep}{0in}
\setlength{\footskip}{5pt}

% sections formatting
\titleformat{\section}{
  \vspace{-2pt}\scshape\raggedright\large
}{}{0em}{}[\color{black}\titlerule\vspace{-5pt}]

% ensure that generate pdf is machine readable/ATS parsable
\pdfgentounicode=1

% custom commands
\newcommand{\cvitem}[1]{
  \item\small{
    {#1\vspace{-2pt}}
  }
}

\newcommand{\cvheading}[4]{
  \vspace{-2pt}\item
    \begin{tabular*}{\textwidth}[t]{l@{\extracolsep{\fill}}r}
      \textbf{#1} & #2 \\
      \small#3 & \small #4 \\
    \end{tabular*}\vspace{-7pt}
}

\newcommand{\cvheadingstart}{\begin{itemize}[leftmargin=0in, label={}]}
\newcommand{\cvheadingend}{\end{itemize}}
\newcommand{\cvitemstart}{\begin{itemize}[label=\textopenbullet]\justifying}
\newcommand{\cvitemend}{\end{itemize}\vspace{-5pt}}

\newcommand{\cvskill}[2]{
  \textcolor{black}{\textbf{#1}}\hfill
  \foreach \x in {1,...,5}{%
    \space{\ifnumgreater{\x}{#2}{\color{black!80!white!20}}{\color{black}}\faSquare}}\par%
  \vspace{-2pt}
}

\renewcommand\labelitemii{$\vcenter{\hbox{\footnotesize$\bullet$}}$}

\begin{document}

% contact information
\begin{center}
  \textbf{\LARGE\scshape 김윤성} \\
  \vspace{1pt}\small
  \href{mailto:p4ranlee@gmail.com}{\textbf{이메일}}
  $\ \diamond\ $ ysk@kzalloc.com
  \textbf{전화}
  $\ \diamond\ $+82 10-3081-9242
  \href{https://github.com/kzall0c}{\textbf{GitHub}}
  \href{https://www.linkedin.com/in/yunseong-kim-linux-kernel/}{\textbf{LinkedIn}}
\end{center}

\section{경력}
\cvheadingstart
  \cvheading
    {에릭슨 LG}{대한민국 서울}
    {시스템 소프트웨어 개발자}{2023.07 - 현재}
  \cvitemstart
    \cvitem{이동성 및 제어 개발팀 (Mobility and Control Development Team)}
  \cvitemend

  \cvheading
    {안랩}{대한민국 분당구}
    {시스템 소프트웨어 개발자}{2022.01 - 2023.06}
  \cvitemstart
    \cvitem{EDR(엔드포인트 탐지 및 대응) 개발팀}
  \cvitemend

  \cvheading
    {이노와이어리스}{대한민국 분당구}
    {소프트웨어 개발자}{2018.08 - 2021.12}
  \cvitemstart
    \cvitem{서버 소프트웨어 개발팀}
\cvitemend

\cvheadingend

\section{학력}
\cvheadingstart
  \cvheading
    {한국항공대학교}{대한민국 고양시}
    {전자 및 항공전자공학 학사}{2010.03 - 2018.02}
  \cvheading
    {리눅스 재단(Linux Foundation)}{온라인}
    {SKF201: 위협 모델링(Threat Modeling)}{2025.09 - 2025.09}
  \cvheading
    {리눅스 재단(Linux Foundation)}{온라인}
    {LFD441: 리눅스 커널과 보안(Security and the Linux Kernel)}{2024.07 - 2024.07}
\cvheadingend

\section{주요 성과}
\cvheadingstart
\item
\cvitemstart
  \cvitem{\textbf{리눅스 커널 보안 개발 (Linux Kernel Security Development)}}
  \begin{itemize}
    \item \textbf{LSM}, \textbf{kprobe}, \textbf{eBPF} 기반의 보안 모듈을 개발
    \item \textbf{perf 도구}를 사용하여 CPU, 메모리, I/O 성능을 프로파일링 및 최적화
  \end{itemize}

\cvitem{\textbf{\href{https://lore.kernel.org/all/?q=Yunseong}{리눅스 커널 업스트림 기여 (Upstream Linux Kernel Contributions)}}}
  \begin{itemize}
    \item 4건의 보안 취약점을 보고하고, 메인라인 커널에 패치 병합:
    \href{https://lore.kernel.org/linux-cve-announce/2025081651-CVE-2025-38510-f67d@gregkh/}{\#1},
    \href{https://lore.kernel.org/linux-cve-announce/2024122721-CVE-2024-53186-7c05@gregkh/}{\#2},
    \href{https://lore.kernel.org/linux-cve-announce/2024080740-CVE-2024-42235-efae@gregkh/}{\#3},
    \href{https://lore.kernel.org/linux-cve-announce/2024072929-CVE-2024-41021-f857@gregkh/}{\#4}
    \item 안정화(LTS) 커널 버전으로 취약점 수정 백포팅 수행
  \end{itemize}

 \cvitem{\href{http://perf.wiki.kernel.org/}{\textbf{Linux perf tools 위키 유지보수자}};}
 {\textbf{uftrace 기여:} \href{https://github.com/namhyung/uftrace/pulls?q=author\%3Aparanlee}{\#1}, \href{https://github.com/namhyung/uftrace/pulls?q=author%3Akzall0c+}{\#2}}

  \cvitem{\textbf{5G 코어 시스템 보안 R\&D}}
  \begin{itemize}
    \item \textbf{eBPF + Rust (aya-rs)} 기반의 NGAP 방화벽 개발
    \item \textbf{ASN.1 코덱}부터 \textbf{5G NGAP 프로토콜 처리}까지 취약점 탐색
  \end{itemize}

  \cvitem{\textbf{클라우드 네이티브 분산 시스템}}
  \begin{itemize}
    \item \textbf{Erlang/OTP} 기반 대규모 분산 5G 애플리케이션 개발
  \end{itemize}

\cvitemend
\cvheadingend

\section{기술 역량}
\cvheadingstart
\item
\cvitemstart
    \cvitem{\textbf{Linux Kernel Security Module} 개발: QEMU, LSM, crash tools, ftrace, perf tools, Wireshark 활용}
    \cvitem{C/C++ 기반의 강력한 프로그래밍 및 디버깅 능력}
    \cvitem{\textbf{Fuzzing}: \textbf{syzkaller}, \textbf{AFL++} 사용 경험; Git 버전 관리, Jenkins CI/CD, Erlang/OTP, Python, shell 스크립팅}
    \cvitem{\textbf{리눅스 커널 CI/CD 및 검증 자동화}}
    \begin{itemize}
      \item \textbf{VMware ESXi, vSphere} 기반 가상화 인프라에서 여러 Jenkins 파이프라인 관리
    \end{itemize}
\cvitemend
\cvheadingend

\section{발표 및 출판}
\cvheadingstart
\item
\cvitemstart

\cvitem{\href{https://youtu.be/__DURK3p218?si=aPn0f2Kixq3V4P7T}{{\textbf{DebConf25:} \textit{Rethinking Cryptography in the Linux Kernel – 포스트 양자 암호(PQC) 시대를 위한 준비}}} }

\cvitem{\href{https://youtu.be/nGxzzORZm1o?si=1k9Mz54uqYNatnO1}{\textbf{Head First Reporting of Linux Kernel CVEs:} \textit{커널 퍼저의 실전 활용}}}
\begin{itemize}
  \item syzkaller를 이용한 KSMBD 및 s390x 취약점 분석 및 리눅스 커널 CVE 발견 사례 발표
\end{itemize}

  \cvitem{\href{https://youtu.be/RRJN8JS5xAE?si=_X-7-OyXzPrcyGes}{{\textbf{DebConf24:}} \textit{perf In Action}} - 페이지 폴트 비교 및 Linux perf를 이용한 퍼징 활용 사례 소개}

  \cvitem{\textbf{\href{https://www.linkedin.com/posts/yunseong-kim-linux-kernel_xilinx-fpga-linux-activity-7341153544160362496-_sBg?utm_source=share&utm_medium=member_desktop&rcm=ACoAADZ7vgEBrd70u0-SnyUUHL-swTgyhcvNQGE}{ZYBO Z7-20 AES 암호화 가속기}} - \textit{\textbf{PetaLinux v2025.1}} 환경에서 HW/SW 공동 개발 수행}

  \cvitem{\href{https://www.hackster.io/p4ranlee/zybo-z7-20-webcam-on-the-petalinux-c28709}{{\textbf{Hackster.io:}} \textit{PetaLinux 프로젝트에서 ZYBO Z7-20 웹캠 구현}}}
  
\cvitemend
\cvheadingend

\end{document}
